\documentclass[11pt]{article}
\usepackage[T1]{fontenc}
\usepackage[utf8]{inputenc}
\usepackage{lmodern}
\usepackage{geometry}
\usepackage{amsmath,amssymb,amsthm,mathtools}
\usepackage{bm}
\usepackage{hyperref}
\usepackage{enumitem}
\usepackage{microtype}
\geometry{margin=1in}
\hypersetup{colorlinks=true,linkcolor=blue,citecolor=blue,urlcolor=blue}
\title{Relational Actualism I:\\Kernel, Closure, and the Engine of Becoming}
\author{Joshua F. Sandeman}
\date{Draft prepared April 2026}

\begin{document}
\maketitle

\begin{abstract}
We present the first paper of a three-paper formulation of Relational Actualism (RA). The aim of this paper is not to catalogue every empirical application of the programme, but to state its kernel cleanly, separate kernel from implementation, and derive the covariant growth architecture that underwrites the rest of the suite. The primitive commitment of RA is that irreversible actualization events are the basic physical reality and that the universe is a growing causal directed acyclic graph rather than a timeless manifold. We organize the framework into five epistemic layers: kernel, closure, implementation, interpretation, and speculation. We then state a five-axiom kernel, the seven-row closure spine that pins the implementation to the Benincasa--Dowker--Glaser structure in four dimensions, and the five-step Engine of Becoming that replaces both the Schrodinger equation and the Einstein field equations as fundamental. The paper also presents the actualization criterion, the Local Ledger Condition, causal invariance, proper time as event count, and the architectural separation between proposal, inscription, and metric update that yields the null Bose--Marletto--Vedral prediction. The result is a complete draft of the conceptual and formal opening paper for the three-paper RA suite.
\end{abstract}

\section{Introduction: The Self-Writing Universe}
Physics has been extraordinarily successful at describing structures, symmetries, and regularities, but it has been markedly less successful at representing the bare fact that the world happens. In both quantum mechanics and general relativity, time enters as part of a formal structure already laid out: a parameter in one case, a coordinate in the other. The equations describe relations among states or fields, but they do not contain, as a primitive, the physical fact that new events come to be.

Relational Actualism begins by refusing to explain happening in terms of something more basic. Happening is primitive. The universe is not a system placed inside a pre-existing arena and made to obey external laws. It is a self-writing causal history. What exists fundamentally is the growing record of irreversible actualization events and the causal relations they inscribe.

This first paper has four jobs. First, it states the kernel of RA in its leanest form. Second, it separates that kernel from the specific 4D implementation used in the current programme. Third, it states the closure results that force the implementation into a tight mathematical spine rather than leaving it as a discretionary modelling choice. Fourth, it presents the Engine of Becoming: the five-step covariant stochastic graph rewriting system from which the later papers derive matter, forces, gravitation, and complexity.

The strategy throughout is epistemically explicit. We distinguish what belongs to the irreducible kernel from what belongs to the current implementation, and we distinguish theorem, published import, exact computation, structural interpretation, and open target. This separation is not cosmetic. It is part of the content of the programme.

\section{Kernel, Closure, and Implementation}
\subsection{The five epistemic layers}
For clarity, we divide the framework into five layers.
\begin{enumerate}[label=\arabic*.]
    \item \textbf{Kernel.} What is philosophically and physically primitive.
    \item \textbf{Closure.} Uniqueness and consistency results that pin the implementation to a narrow mathematical form.
    \item \textbf{Implementation.} The concrete 4D theory used in the present RA suite.
    \item \textbf{Interpretation.} The physical reading of the mathematical structures.
    \item \textbf{Speculation.} Conjectural extensions not yet forced by the established structure.
\end{enumerate}

This architecture prevents a recurring failure mode in ambitious foundational programmes: letting the kernel inherit every strength of the implementation while simultaneously letting the implementation borrow every necessity of the kernel. RA is strongest when the distinction is explicit.

\subsection{The five kernel axioms}
The irreducible kernel may be stated as follows.

\paragraph{Axiom A: Event primacy.} The primitive physical reality is not a field on a manifold but a growing set of events related by causal precedence.

\paragraph{Axiom B: Irreversible inscription.} Some interactions are genuinely irreversible and write permanent causal facts.

\paragraph{Axiom C: Open future.} At any stage there is a structured adjacent possible: multiple admissible continuations weighted by a physical amplitude assignment.

\paragraph{Axiom D: Nontrivial actuality threshold.} Not every possibility becomes actual. There is a physically meaningful criterion that distinguishes merely virtual continuations from actualized events.

\paragraph{Axiom E: Stable motifs.} Persistent entities are not enduring substances but self-reproducing motifs in the growth dynamics.

These axioms do not by themselves specify the exact local action, the dimension, or the detailed numerical content of the theory. They specify what kind of theory RA is.

\section{The Seven-Row Closure Spine}
The current RA implementation is not presented as an arbitrary elaboration placed on top of the kernel. It is constrained by a seven-row closure spine:
\begin{center}
\begin{tabular}{p{0.08\textwidth} p{0.34\textwidth} p{0.40\textwidth}}
1 & A local combinatorial action exists & In 4D the BDG action is uniquely selected as the relevant discrete curvature action.\\
2 & Conservation is local & The Local Ledger Condition holds at vertices.\\
3 & Amplitude depends on causal past & Local amplitude depends only on the causal past of the candidate event.\\
4 & Spacelike ordering is irrelevant & Causal invariance holds: spacelike reorderings do not change physical amplitudes.\\
5 & Dense coarse-graining yields a continuum & The BDG dense limit recovers Einstein--Hilbert structure.\\
6 & Actualization has a nontrivial threshold & Positive BDG score implements the discrete actualization filter.\\
7 & Stable patterns define matter & Stable motif classes close under growth and provide the matter taxonomy.\\
\end{tabular}
\end{center}

The philosophical consequence is important. The gap between kernel and implementation is not filled by taste. It is filled, as far as the present programme goes, by a mixture of uniqueness results, formal proofs, and exact computations.

\section{The Growing Causal Graph}
Let $G_n=(V_n,E_n)$ denote the actualized universe after $n$ events, where $V_n$ is the set of actualized vertices and $E_n$ is the set of directed causal edges. The graph is acyclic. An edge $u\to v$ means that $u$ is in the causal past of $v$. The acyclicity is not a convenience but the discrete expression of irreversibility.

The present boundary of the universe is the set of vertices with no actualized future successors. The future is not a pre-existing region of a block manifold. It is the structured set of candidate extensions of the current graph.

\subsection{Proper time as count}
If $\gamma$ is a worldline through the graph, then the primitive definition of proper time is
\begin{equation}
\tau(\gamma)=N(\gamma)\, t_P,
\end{equation}
where $N(\gamma)$ is the number of actualization events along the worldline. The continuum proper-time integral is a large-$N$ approximation to this integer count, not the other way around.

\subsection{The local action}
For the current 4D implementation the local discrete curvature score is the Benincasa--Dowker--Glaser expression
\begin{equation}
S_{\mathrm{BDG}}(v)=1-N_1(v)+9N_2(v)-16N_3(v)+8N_4(v),
\end{equation}
up to the conventional Planck-scale normalization. Here $N_k(v)$ counts inclusive order-interval structures in the causal past of $v$ of depth $k$ in the 4D BDG scheme.

\section{The Actualization Criterion}
RA requires a criterion separating unrealized possibility from realized event. In the continuum and AQFT language, the fundamental statement is that actualization occurs when an interaction produces an irreversible increase in quantum relative entropy relative to the vacuum reference state:
\begin{equation}
\Delta S(\rho\Vert\sigma_0)>0.
\end{equation}

In the discrete implementation, this is represented by the positive BDG filter:
\begin{equation}
S_{\mathrm{BDG}}(v)>0.
\end{equation}
The point is not that the entropic and graph-theoretic criteria are two unrelated proposals. The claim of the current implementation is that the second is the discrete physical realization of the first.

This hierarchy is best stated in four levels.
\begin{enumerate}[label=\roman*.]
    \item The deepest criterion is record-forming irreversibility.
    \item In AQFT language this is captured by positive entropic increment.
    \item In perturbative QFT, on-shell emission is a sufficient operational implementation in controlled regimes.
    \item In the discrete causal graph, positive BDG score is the local filter.
\end{enumerate}

\section{The Local Ledger Condition}
The local conservation principle of the graph is the Local Ledger Condition (LLC). If $q(v)$ denotes the charge vector carried by a vertex, then for every actualization vertex $v$,
\begin{equation}
\sum_{u:(u,v)\in E} q(u\to v)=\sum_{w:(v,w)\in E} q(v\to w).
\end{equation}

The LLC is not an added bookkeeping rule imposed on an otherwise complete theory. It is the discrete conservation structure without which a growing causal ledger would not cohere as a physical universe. In the continuum limit the familiar conservation equations arise as coarse-grained expressions of the same local balance.

\section{The Engine of Becoming}
The core dynamics of RA is a five-step process carried out at each stage of graph growth.

\subsection*{Step 1: Survey}
Identify the causal frontier and the set of candidate continuations compatible with the existing partial order.

\subsection*{Step 2: Propose}
Assign complex amplitudes to candidate extensions from the full set of compatible causal histories. This is the layer of structured possibility, not yet actuality.

\subsection*{Step 3: Filter}
Discard candidates that fail the actuality criterion. In the discrete implementation this is the BDG positivity condition.

\subsection*{Step 4: Inscribe}
Select one admissible candidate according to the quantum measure and write it irreversibly into the graph.

\subsection*{Step 5: Update}
Recompute the macroscopic description from the enlarged actualized graph. Geometry belongs here, not at Step 2.

This separation is architecturally decisive. Proposal and geometry update are different stages. The weighted space of unrealized continuations does not itself source the metric. The metric is sourced only by what has actually been inscribed.

\section{Causal Invariance}
Any sequential growth dynamics for a relativistic theory must avoid dependence on a preferred foliation. The discrete RA condition is causal invariance: if two candidate events are spacelike relative to one another, then the total amplitude assigned to the resulting physical history is independent of the order in which the bookkeeping sequence lists them.

In symbolic form, if $v_A$ and $v_B$ are spacelike and admissible, then the physical quantum measure for the completed history satisfies
\begin{equation}
\mu(v_A\ \text{then}\ v_B)=\mu(v_B\ \text{then}\ v_A).
\end{equation}
The combinatorial core of this statement is that local amplitudes depend only on causal pasts and hence commute under transpositions of spacelike-separated insertions.

Causal invariance is the discrete analogue of Lorentz covariance. It is not covariance of a background manifold. It is covariance of the growth law itself.

\section{Proper Time, Energy, and Causal Bandwidth}
Once proper time is treated as event count, the identification between energy and actualization rate becomes natural. For a stable motif of rest mass $m$, the rest-frame actualization frequency is interpreted by the usual bridge relation
\begin{equation}
f_0=\frac{mc^2}{h}.
\end{equation}
This should not be read as importing continuum quantum mechanics as primitive. It should be read as the late translation between human units and event-count dynamics.

The speed of light is then interpreted as the maximal causal propagation rate,
\begin{equation}
c=\frac{\ell_P}{t_P},
\end{equation}
with massless propagation corresponding to full causal bandwidth devoted to spatial propagation and massive propagation corresponding to bandwidth partly consumed by motif maintenance.

\section{The BMV Null Prediction and the Step-2/Step-5 Separation}
The Bose--Marletto--Vedral protocol is the clean near-term experimental consequence of the architecture. In the standard quantum-gravity reading, branch-dependent Newtonian phases can entangle two masses because each branch of the matter superposition sources its own gravitational field. In RA this is denied at the level of architecture.

The metric belongs to Step 5. Superposed candidate branches belong to Step 2. Therefore branch-coherent possibility does not by itself generate a branch-resolved metric. The theory predicts no gravity-mediated entanglement in the branch-coherent pre-actualization regime. A confirmed positive BMV signal would therefore strike directly at the Step-2/Step-5 separation.

\section{What Paper I Establishes}
This first paper establishes the conceptual and formal entry point for the three-paper RA suite.
\begin{enumerate}[label=\arabic*.]
    \item The universe is modeled as a growing causal ledger rather than a static block manifold.
    \item The kernel of the theory is compact: event primacy, irreversible inscription, open future, actuality threshold, and stable motifs.
    \item The current 4D implementation is tied to the kernel by a closure spine rather than by discretionary modelling.
    \item The Engine of Becoming provides a covariant growth law with a sharp distinction between possibility, actuality, and post-actualization geometry.
    \item Proper time, conservation, and covariance are recast in discrete causal terms.
\end{enumerate}

The second and third papers take these ingredients forward into the matter/force sector and the gravitation/cosmology/complexity sector respectively.

\section{Conclusion}
Relational Actualism begins not with fields, particles, or manifolds, but with the physically primitive fact that irreversible events occur and accumulate. The universe is a growing causal ledger. Geometry is its large-scale order. Matter is the persistence of stable motifs under the growth rule. Time is the fact that the ledger grows.

The purpose of this first paper has been to state that architecture cleanly. A foundational programme becomes stronger, not weaker, when it separates kernel from implementation and closure from interpretation. On that architecture the rest of the RA suite can proceed in an orderly way: first to matter, forces, and particle persistence; then to gravitation, cosmology, and complexity.

\begin{thebibliography}{99}
\bibitem{BD2010} D. M. T. Benincasa and F. Dowker, ``The scalar curvature of a causal set,'' \emph{Physical Review Letters} \textbf{104} (2010), 181301.
\bibitem{Dowker2013} F. Dowker, ``Introduction to causal sets and their phenomenology,'' \emph{General Relativity and Gravitation} \textbf{45} (2013), 1651--1667.
\bibitem{Haag1996} R. Haag, \emph{Local Quantum Physics}, 2nd ed., Springer, 1996.
\bibitem{Lovelock1971} D. Lovelock, ``The Einstein tensor and its generalizations,'' \emph{Journal of Mathematical Physics} \textbf{12} (1971), 498--501.
\bibitem{RideoutSorkin2000} D. Rideout and R. D. Sorkin, ``A classical sequential growth dynamics for causal sets,'' \emph{Physical Review D} \textbf{61} (2000), 024002.
\bibitem{Rovelli1996} C. Rovelli, ``Relational quantum mechanics,'' \emph{International Journal of Theoretical Physics} \textbf{35} (1996), 1637--1678.
\bibitem{Sorkin1994} R. D. Sorkin, ``Quantum mechanics as quantum measure theory,'' \emph{Modern Physics Letters A} \textbf{9} (1994), 3119--3127.
\bibitem{SandemanRAQM} J. F. Sandeman, ``The Relational Universe: Interaction and the Transition of the Possible to the Actual,'' preprint, 2026.
\bibitem{SandemanRAGC} J. F. Sandeman, ``Relational Actualism Gravity and Cosmology,'' preprint, 2026.
\bibitem{SandemanRAEB} J. F. Sandeman, ``Relational Actualism: The Engine of Becoming,'' preprint, 2026.
\end{thebibliography}

\end{document}
